\section{Data Quality Control\label{sec: Data Quality Control}}
%\subsection{Data Annotation of Subgoals: %\label{sec:Data Annotation of sub-goals} }
%To get more detailed progress, we annotate subgoals for these tasks. The details are listed below:

%\paragraph{For BabyAI, }
%our annotation process involves the use of an interactive interface for identifying critical observations, as illustrated in Figure \ref{fig: sub-goal annotations 1}. Recognizing that some goals were impractical for Language Learning Models (LLMs) due to the lack of consistent examples of completed tasks, we revise the goal description to opt for more feasible objectives. Furthermore, we refine our task descriptions to ensure that the labeled observations are necessary, such as specifying the tools used for task completion.
%\paragraph{For Jericho, }
%our annotation process involves the use of an interactive interface for identifying critical observations, as illustrated in Figure \ref{fig: sub-goal annotations 1}. Recognizing that some goals were impractical for Language Learning Models (LLMs) due to the lack of consistent examples of completed tasks, we revise the goal description to opt for more feasible objectives. Furthermore, we refine our task descriptions to ensure that the labeled observations are necessary, such as specifying the tools used for task completion.
%\paragraph{For PDDL,}
%our annotation process involves the use of an interactive interface for identifying critical observations, as illustrated in Figure \ref{fig: sub-goal annotations 1}. Recognizing that some goals were impractical for Language Learning Models (LLMs) due to the lack of consistent examples of completed tasks, we revise the goal description to opt for more feasible objectives. Furthermore, we refine our task descriptions to ensure that the labeled observations are necessary, such as specifying the tools used for task completion.
%\paragraph{For Game24, }
%our annotation process involves the use of an interactive interface for identifying critical observations, as illustrated in Figure \ref{fig: sub-goal annotations 1}. Recognizing that some goals were impractical for Language Learning Models (LLMs) due to the lack of consistent examples of completed tasks, we revise the goal description to opt for more feasible objectives. Furthermore, we refine our task descriptions to ensure that the labeled observations are necessary, such as specifying the tools used for task completion.
%\paragraph{For Webshop, }
%our annotation process involves the use of an interactive interface for identifying critical observations, as illustrated in Figure \ref{fig: sub-goal annotations 1}. Recognizing that some goals were impractical for Language Learning Models (LLMs) due to the lack of consistent examples of completed tasks, we revise the goal description to opt for more feasible objectives. Furthermore, we refine our task descriptions to ensure that the labeled observations are necessary, such as specifying the tools used for task completion.
%\paragraph{For Webarena, }
%our annotation process involves the use of an interactive interface for identifying critical observations, as illustrated in Figure \ref{fig: sub-goal annotations 1}. Recognizing that some goals were impractical for Language Learning Models (LLMs) due to the lack of consistent examples of completed tasks, we revise the goal description to opt for more feasible objectives. Furthermore, we refine our task descriptions to ensure that the labeled observations are necessary, such as specifying the tools used for task completion.
%\paragraph{For Tool, }
%We first meticulously design goals for each task.
%In formulating each goal, we ensure that it requires multiple steps to accomplish and manually annotate a golden action path to achieve it.
%Following the golden action path, we collect the corresponding observations, which we post-process as subgoals.
%We guarantee the necessity of each action in the golden action path, thus, the obtained observations are essential, making each subgoal a requisite for accomplishing the goal.
%However, for the Sheet task, since a unique set of subgoals could not be determined, we do not annotate subgoals for Sheet task.
% Instead, we adopt evaluation metrics similar to those used for web tools, calculating a matching score for spreadsheet after each action step in relation to the final spreadsheet.
\begin{figure*}[t!]
\centering 
\includegraphics[width=\textwidth]{figures/check.pdf} 
\caption{An illustration of our sub-goal checking interface. We develop an interactive interface for annotators to checking sub-goals. Firstly, annotators play and pass the game with the interface. The reward score for each step will be given based on the labeled score. If the annotators are dissatisfied with the reward, annotators will record them and the corresponded environment will be discussed by more annotators and annotated again."}
\label{fig: sub-goal annotations 2} 
\end{figure*}




\begin{table}[!t]
\centering
\resizebox{0.6\linewidth}{!}{\begin{tabular}{cc@{}}
\toprule
\textbf{Task} & \textbf{Proportion of environments with annotation errors} \\ \midrule
 AlfWorld & 10.0\%   \\
 ScienceWorld & 0\%   \\
 Babyai & 4.2\%    \\
Jericho  & 25\%  \\
PDDL & 5\%   \\
 WebShop &  0\% \\
 WebArena &   0\% \\
 Tool-Query &  0\%  \\
 Tool-Operation  &  0\% \\ \bottomrule
\end{tabular}}
\caption{The proportion of environments with annotation errors in the second round of data checking. Environments identified with errors are subsequently analyzed to determine the underlying causes, and any environments exhibiting similar errors are amended collectively.}
\label{tab: ration of environments with annotation erros}
\end{table}
To ensure the quality of labeled sub-goals, we conducted three rounds of data verification for each labeled sub-goal. We developed an interactive interface through which inspectors complete tasks and observe the reward scores obtained at each step. If the inspector deems the reward score assigned during interaction with an environment to be unreasonable, additional annotators will engage in a discussion to determine if modifications to the labeled sub-goals are necessary.

The first round of verification is a self-check. Each annotator is required to carefully review the labeled tasks in every environment they are responsible for. The second round involves a sampled inspection by two annotators for each task. They examine a sample of 5-10 items from different sub-tasks within the task and document the proportion of issues identified, as presented in Table \ref{tab: ration of environments with annotation erros}. The third round is conducted by an annotator who is well-acquainted with the various tasks, who then performs a sampled review of all tasks.
% \subsection{Examples of goals and trajectory}

% % \renewcommand{\arraystretch}{1.2} % 1.2倍行高
\begin{table*}[t!]
\centering
\scriptsize
\caption{Examples of goals for the 4 task categories in \BenchName{}, along with a sampled step of the trajectory and progress rate.
The trajectory is generated by GPT-4.
Some lengthy observations are omitted with ``$\ldots$'' for brevity. The task name in the table uses an abbreviation, the full name can be found in \S \ref{section: task decomposition}.
}
\vspace{-2mm}
\resizebox{\textwidth}{!}{
\begin{tabular}{p{0.3cm}|p{12.50cm}}
\toprule
\multicolumn{1}{c|}{\textbf{Task}}  & \multicolumn{1}{c}{\textbf{Goal \& Trajectory}} \\
\midrule
\midrule
\multirow{2}[5]{*}{\textbf{ALF}} 
&
\textbf{Goal: put a clean egg in microwave.}
% id: 23
\\
&
\stepbox{Step 02}
\hspace{0.18cm}\tcv{Action:} open fridge 1
\ \tco{Observation:} You open the fridge 1. The fridge 1 is open. In it, you see a apple 2, a egg 1, a lettuce 1, a pan 2, a plate 1, and a tomato 1.
\hfill \stepbox{Progress Rate: 0.00 $\rightarrow$ 0.25}
% 218, 238, 253
% 183, 220, 252  0.71, 0.86, 0.99
% \\
% \midrule
% \midrule
% \multirow{2}[5]{*}{\textbf{SW}} 
% &
% \textbf{Goal: Your task is to boil apple juice until it reaches a gaseous state.}
% % id: 36
% \\
% &
% \stepbox{Step 25}
% \hspace{0.05cm}\tcv{Action:} use thermometer in inventory on apple juice in metal pot
% \ \tco{Observation:} the thermometer measures a temperature of 23 degrees celsius
% \hfill \stepbox{Progress Rate: 0.75 $\rightarrow$ 0.88}
% \\
% \midrule
% \midrule
% \multirow{2}[5]{*}{\textbf{BA}} 
% &
% \textbf{Goal: pick up the box}
% % id: 59
% \\
% &
% \stepbox{Step 07}
% \hspace{0.05cm}\tcv{Action:} toggle and go through green closed door 1
% \quad\tco{Observation:}  In front of you in this room, you can see several objects: There is a red box 1 $\ldots$ 
% \hfill \stepbox{Progress Rate: 0.33 $\rightarrow$ 0.66}
\\
\midrule
\midrule
\multirow{2}[5]{*}{\textbf{JC}} 
&
\textbf{Goal: Get out of the house. Then escape the city without getting caught via driving.}
% id: 0
\\
&
\stepbox{Step 29}
\hspace{0.18cm}\tcv{Action:} take a shower
\ \tco{Observation:} You step into the shower, turn on the water, and within a few moments you feel like a new man. But no time to dawdle - you hop back out again and dry off in record time $\ldots$ \hfill \stepbox{Progress Rate: 0.43 $\rightarrow$ 0.57}
\\
\midrule
\midrule
% \multirow{2}[5]{*}{\textbf{PL}}
% &
% \textbf{Goal: The goal is to satisfy the following conditions: Wheel r1 is inflated., r1 is on the-hub1., w1 is in boot.}
% % id: 50
% \\
% &
% \stepbox{Step 08}
% \hspace{0.18cm}\tcv{Action:} Put-away w1 in boot.
% \ \tco{Observation:} W1 is in boot.
% \hfill \stepbox{Progress Rate: 0.00 $\rightarrow$ 0.33}
% \\
% \midrule
% \midrule
% \multirow{2}[5]{*}{\textbf{WS}} 
% &
% \textbf{Goal: i want a noise cancelling cosycost usb microphone, and price lower than 70.00 dollars}
% % id: 1
% \\
% &
% \stepbox{Step 02}
% \hspace{0.18cm}\tcv{Action:} search[noise cancelling cosycost usb microphone]
% \ \tco{Observation:} WEB PAGE: \{ [Back to Search] Page 1 (Total results: 15) [Next $>$] $\ldots$
% \hfill \stepbox{Progress Rate: 0.00 $\rightarrow$ 0.33}
% \\
% \midrule
% \midrule
\multirow{2}[5]{*}{\textbf{WA}} 
&
\textbf{Goal: Display the list of issues in the kkroening/ffmpeg-python repository that have labels related to questions}
% id: 1
\\
&
\stepbox{Step 05}
\hspace{0.18cm}\tcv{Action:} click [5398]
\ \tco{Observation:} Tab 0 (current): Issues · Karl Kroening / ffmpeg-python · GitLab [6573] RootWebArea `Issues · Karl Kroening / ffmpeg-python · GitLab' focused: True [6620] link $\ldots$
\hfill \stepbox{Progress Rate: 0.25 $\rightarrow$ 0.50}
% \\
% \midrule
% \midrule
% \multirow{2}[5]{*}{\textbf{TQ}} 
% &
% \textbf{Goal: Avatar versus Forrest Gump, which has a higher rating?}
% % id: 22
% \\
% &
% \stepbox{Step 04}
% \hspace{0.18cm}\tcv{Action:} get\_movie\_details with Action Input: \{``movie\_id": ``13"\}
% \ \tco{Observation:} \{`title': `Forrest Gump', `vote\_average': 8.476, `budget': 55000000, $\ldots$ 
% \hfill \stepbox{Progress Rate: 0.60 $\rightarrow$ 0.80}
\\
\midrule
\midrule
\multirow{2}[7]{*}{\textbf{TO}} 
&
\textbf{Goal: In ``Sheet17'', calculate and complete the ``Profit'' of the products in the table based on the sales information of the products. And then, sort the table in descending order by ``Profit''.}
% id: 36
\\
&
\stepbox{Step 07}
\hspace{0.1cm}\tcv{Action:} update\_cell\_by\_formula with Action Input: \{``operator'': ``PRODUCT'',  `` start\_position'': ``C8'', ``end\_position'': ``D8'', ``result\_position'': ``E8''\}
\ \tco{Observation:}
[[`Product', `Category',
%`Price',
%`Sold out', 
% `Profit'], [`banana', `Fruits', 
$\ldots$ 
\hfill \stepbox{Progress Rate: 0.47 $\rightarrow$ 0.49}
\\
\bottomrule
\end{tabular}}

\label{tab:goal_trajectory}
% \vspace{-5mm}
\end{table*}

% \newpage